\section{Pre-registered design}

The evaluation specification was frozen on 2026-07-06, before any simulator
code existed, and committed to an append-only repository; the frozen text
with its complete amendment log is released with this paper \citep{evalspec2026}.
This section restates its commitments; the specification governs where they
differ.

Primary metric. Task-level goodput per provisioned GPU-hour: value-weighted
completed flows per hour, normalized by provisioned capacity. The
normalization is the honesty mechanism: capacity a reservation holds idle
lowers the discipline's own score arithmetically, so stranding, the
mechanism's known cost, is priced inside the headline number rather than
negotiated beside it. Value weights are fixed in the specification and
varied only in a declared sensitivity sweep. Related serving metrics score
completed requests against latency targets \citep{mars2026}; ours differs in unit
and denominator, flow completions valued per provisioned GPU-hour, so that
held-but-idle guarantees are charged to the discipline that holds them.

Thresholds, fixed before data. Support required a median improvement of at
least 10 percent over the best baseline, selected per cell as the strongest
tuned opponent in that cell (deliberately conservative against the
mechanism), across at least 60 percent of the core region, with
interactive-traffic p95 time-to-first-token degradation held under 5
percent. Below 5 percent median improvement, or on guardrail violation
across the region: kill. Between: a weak result. The verdict is computed
mechanically by released code as a pure function of the run manifests.

Region. The core region spans interactive-to-agentic traffic mixes (90/10,
70/30, 50/50 by token volume), offered loads 0.7 to 1.5 times measured
capacity, idle-gap distributions with medians of 20, 60, and 300 seconds,
and cluster sizes of 8 to 32 devices (Amendment 3, below). Corner cells
were pre-declared excluded; none required exclusion in the final grid.

Baselines and fairness. Six baselines: vLLM-style priority scheduling
with preemption and paged cache \citep{vllm2023}; FastServe-style multi-level
feedback queues \citep{wu2023fastserve}; Niyama-style deadline-aware class
scheduling, implemented here with a deterministic demotion ladder
\citep{niyama2025}; CONCUR-style AIMD admission with pause/resume
\citep{concur2026}; a MARS-style external admission controller with
cost-benefit cache retention \citep{mars2026}, added by Amendment 1; and FCFS
continuous batching as the no-control floor. Every tunable policy,
including ours, received an identical tuning budget: an 8-point
pre-registered grid, tuned per workload family on validation seeds and
scored on disjoint held-out seeds, with a build-failing test asserting
no policy's grid is larger than any other's. Each baseline is a
policy-family abstraction rather than a reimplementation of any single
system: the families are named for the work that characterises them, and
the comparison is against the scheduling discipline each represents. The
tuned parameters and both scores are released (Table~\ref{tab:roster}, in Appendix G).

The registration carries three numbered amendments and one integrity
ruling, each dated, reasoned, and adopted before the final grid ran: a
baseline added because a negative headline makes baseline completeness
load-bearing, a measurement horizon corrected, and the grid's device axis
reduced after measured infeasibility. All were adopted under a one-fix
clause barring any further measurement correction in either direction,
which is why the verdict below stands as computed; Appendix E gives the
log in full.
